\centerline{\bf \Large  9 класс}

\problem{\textbf{1}} 
Пюрвя загадал число, которое имеет $2025$ различных делителей, отличных от $1$ и самого числа. Сколько простых делителей может быть у числа Пюрви?

\textbf{Решение:} Заметим, что у числа всего 2027 делителей, считая 1 и само число. Вспомним, как считается количество делителей числа:
если степени вхождения простых чисел в его разложение равны $a_1, a_2, ..., a_n$, то количество делителей равно
$(a_1+1)(a_2+1)...(a_n+1)=2027$.

Но заметим, что 2027 простое число, а значит, в произведении участвует лишь один множитель, то есть число является 2026-ой степенью некоторого простого числа, откуда делаем вывод, что число может иметь только 1 простой делитель.


\textbf{Критерии.} \\
\textit{7 баллов} $-$ Полностью верное решение\\
\textit{0-6 баллов} $-$ Полностью верное решение, но учитывается само число и единица (делителей 2025)\\
\textit{0-6 баллов} $-$ Все остальные продвижения\\
\textit{0 баллов} $-$ Любой ответ без объяснений\\


\problem{\textbf{2}} Решите задачу:

а) Батыр и Очир ходили вокруг школы, оба в одном направлении и каждый с постоянной скоростью. Батыр проходил один круг за 5 минут, а Очир, который был медленнее Батыра, - за некоторое целое число минут. При этом время между встречами тоже равнялось некоторому целому числу минут, причём оно было больше 12. За какое время Очир проходил полный круг?

б) Составьте и решите обратную задачу по схеме: ?; 6; 12.

\textbf{Решение А:} Обозначим через $n$ время в минутах, за которое проходит круг Очир. Тогда $n>5, n \in \mathbb{N}$. Скорость, с которой более быстрый мальчик догоняет медленного равна

$$
\frac{1}{5}-\frac{1}{n}=\frac{n-5}{5 n}\left(\frac{\kappa p y r}{min}\right),
$$


поэтому время между встречами составляет

$$
\frac{5 n}{n-5}=\frac{5(n-5)+25}{n-5}=5+\frac{25}{n-5}(\text { мин }),
$$


а поскольку это целое число, то $(n-5)$ делит 25. Натуральные делители 25 - числа $1 , 5$ и 25 . Тогда $n=6$ или $n=10$ или $n=30$, что соответствует времени между встречами 30 минут, 10 минут и 6 минут. Поскольку время между встречами больше 12 минут, заданному условию удовлетворяет только $n=6$.


\textbf{Примерное условие Б:} Батыр и Очир ходили вокруг школы, оба в одном направлении и каждый с постоянной скоростью. Очир проходил один круг за 6 минут, а Батыр, который был быстрее Очира, - за некоторое целое число минут. При этом время между встречами тоже равнялось некоторому целому числу минут, причём оно было больше 12. За какое время Очир проходил полный круг?

\textbf{Решение Б:} Обозначим через $n$ время в минутах, за которое проходит круг Батыр. Тогда $n<6, n \in \mathbb{N}$. Скорость, с которой более быстрый мальчик догоняет медленного равна

$$
\frac{1}{n}-\frac{1}{6}=\frac{6-n}{6 n}\left(\frac{\kappa p y r}{min}\right),
$$


поэтому время между встречами составляет

$$
\frac{6 n}{6-n}=\frac{36-(6-n)6}{6-n}=\frac{36}{6-n}-6(\text { мин }),
$$


а поскольку это целое число, то $(6-n)$ делит 36 и $\dfrac{36}{6-n}-6>0$. Тогда $n=5$ или $n=4$ или $n=3$ или $n=2$, что соответствует времени между встречами 30 минут, 12 минут, 6 минут и 3 минуты. Поскольку время между встречами больше 12 минут, заданному условию удовлетворяет только $n=5$.

\textbf{Критерии.} \\
\textit{0-4 баллов} $-$ Решение пункта А\\
\textit{0-1 балла} $-$ Составление условия пункта Б \\
\textit{0-2 баллов} $-$ Решение пункта Б\\
\textit{0 баллов} $-$ Любой ответ без объяснений\\


\problem{\textbf{3}} Прямая проходит через точку с координатами $(2025 ; 0)$ и пересекает параболу $y=$ $x^2$ в точках с абсциссами $a$ и $b$. Найдите $\dfrac{1}{a}+\dfrac{1}{b}$.

\textbf{Решение:} Точки $a, b$, являются решениями уравнения $x^2=k(x-2025)$, где $\mathrm{k}-$ коэффициент наклона прямой. Тогда по теореме Виета $a+b=k, a b=2025 k$. Следовательно, $\dfrac{1}{a}+\dfrac{1}{b}=\dfrac{a+b}{a b}=\dfrac{k}{2025 k}=\dfrac{1}{2025}$.


\textbf{Критерии.} \\
\textit{7 баллов} $-$ Полностью верное решение\\
\textit{0-6 баллов} $-$ Все остальные продвижения\\
\textit{0 баллов} $-$ Любой ответ без объяснений\\


\problem{\textbf{4}} 
Дан вписанный четырёхугольник $A B C D$. Прямые $A B$ и $D C$ пересекаются в точке $E$, а прямые $B C$ и $A D$ - в точке $F$. В треугольнике $A E D$ отмечен центр вписанной окружности $I$, а из точки $F$ проведен луч, перпендикулярный биссектрисе угла $A I D$. В каком отношении этот луч делит угол $A F B$ ?

\textbf{Решение:} Заметим, что угол между биссектрисами углов $A E D$ и $A F B$ равен полусумме углов $F A E$ и $F C E$, т.е. $90^{\circ}$. Поэтому угол между биссектрисой угла $A F B$ и лучом $F K$, где $K$ - проекция $F$ на биссектрису угла $A I D$, равен $180^{\circ}-\angle E I K=$ $180^{\circ}-\left(90^{\circ}+\angle A / 2\right)-\left(180^{\circ}-\angle A / 2-\angle D / 2\right) / 2=(\angle D-\angle A) / 4=\angle A F B / 4$ (рис.2), следовательно, и $\angle A F K=\angle A F B / 4$.

\textbf{Критерии.} \\
\textit{7 баллов} $-$ Полностью верное решение\\
\textit{0-6 баллов} $-$ Все остальные продвижения\\
\textit{0 баллов} $-$ Любой ответ без объяснений\\

\problem{\textbf{5}}
На координатной плоскости назовём \textit{бриллиантовым расстоянием} между двумя точками $A$ и $B$, находящимися в узлах, количество клеток, которые пересекает отрезок $AB$. Также назовём \textit{бриллиантовым кругом} радиуса $n$ геометрическое место точек, удалённых от данной на бриллиантовое расстояние $n$. Докажите, что бриллиантовый круг радиуса $n$ с центром в начале координат совпадает со множеством узлов сетки, заданных уравнением $|x| + |y| = n+1$, где $x,y$ - целые тогда и только тогда, когда $n+1$ простое.

\textbf{Решение:}  \textit{Лемма.}
    Рассмотрим клетчатый прямоугольник $m \times n$. Тогда его диагональ пересечёт ровно $m+n - (m, n)$ клеток.
    
    \textit{Доказательство леммы.} Пусть сперва $(m, n) = 1$. Нетрудно заметить, что диагональ пересекает $m - 1$ внутренних горизонталей и $n - 1$ и внутренних вертикалей. Очевидно, что диагональ не может пересечь одновременно внутреннюю вертикаль и горизонталь, иначе это бы означало, что диагональ пришла в узел ещё до того как дошла до правой верхней вершины прямоугольника, что незамедлительно означает, что $(m, n)>1$. Стало быть пересечений диагонали с вертикалями и горизонталями ровно $m+n-2$ штуки. Каждое такое пересечение означает "переход"\, из одной клетки в другую, а значит пересечений с клетками будет на 1 больше. Теперь, если $(m, n) = d$ разобьём его каждую сторону на $d$ равных частей. Тогда диагональ будет проходить через вершины $d$ прямоугольников размером $\frac{m}{d} \times \frac{n}{d}$. В каждом из них она пересекает ровно $\frac{m}{d} + \frac{n}{d} - 1$ клеток. А значит, во всём прямоугольнике пересекает $(\frac{m}{d} + \frac{n}{d} - 1)d = m + n - d$ клеток.

    Пусть $n+1$ - простое. Тогда, по нашей лемме, в любом прямоугольнике $t \times (n + 1 - t)$ диагональ будет пересекать ровно $n+1$ вершину, так как $(t, n+1 - t) = (t, n+1) = 1$. А множество концов таких диагоналей и есть множество вершин, заданных уравнением $|x| + |y| = n + 1$. Отсюда же, на самом-то деле следует, что если бриллиантовый круг совпал с этим множеством, то $n+1$ - простое.

    \textbf{Критерии.} \\
\textit{7 баллов} $-$ Полностью верное решение\\
\textit{0-6 баллов} $-$ Все остальные продвижения\\
\textit{0 баллов} $-$ Любой ответ без объяснений\\